\subsection{Concatenation and finite fibers}
\label{sec:poisson_word_append}

For an arbitrary alphabet $I$, let $\mathcal W$ be its finite words.
The following statements do not require $I$ to be finite.

\begin{definition}[Word concatenation]
\label{def:poisson_word_append}
\lean{PoissonWord.append, PoissonWord.length_append,
PoissonWord.nil_append, PoissonWord.append_nil, PoissonWord.append_assoc}
\leanok
\uses{def:poisson_word_measure}
For words $u$ and $v$, their concatenation $uv$ places all letters of $v$
after all letters of $u$.  Thus $|uv|=|u|+|v|$.
Concatenation is associative, and the empty word is a two-sided identity.
\end{definition}

\begin{definition}[Prefixes and suffixes]
\label{def:poisson_word_split}
\lean{PoissonWord.take, PoissonWord.drop, PoissonWord.length_take,
PoissonWord.length_drop, PoissonWord.length_take_of_le,
PoissonWord.ofFn_take, PoissonWord.ofFn_drop}
\leanok
\uses{def:poisson_word_measure}
For a word $w$ and $k\geq0$, write $w_{<k}$ for its first
$\min(k,|w|)$ letters, and $w_{\geq k}$ for the remaining
$\max(|w|-k,0)$ letters.  In particular,
$|w_{<k}|=k$ whenever $k\leq|w|$.
\end{definition}

\begin{lemma}[Fibers of concatenation]
\label{lem:poisson_word_append_fiber}
\lean{PoissonWord.append_take_drop, PoissonWord.take_append,
PoissonWord.drop_append, PoissonWord.appendFiber,
PoissonWord.appendFiberEquiv, PoissonWord.appendFiberEquiv_apply,
PoissonWord.appendFiberEquiv_symm_apply, PoissonWord.instFintypeAppendFiber,
PoissonWord.card_appendFiber, PoissonWord.finite_append_preimage_singleton,
PoissonWord.ncard_append_preimage_singleton, PoissonWord.append_eq_iff}
\leanok
\uses{def:poisson_word_append, def:poisson_word_split}
For every word $w$ of length $m$, the fiber
\begin{align}
 \{(u,v)\in\mathcal W\times\mathcal W:uv=w\}
 \label{eq:poisson_word_append_fiber}
\end{align}
has exactly $m+1$ elements.  They are parametrized bijectively by
$k\in\{0,\ldots,m\}$ through
\begin{align}
 k\longmapsto (w_{<k},w_{\geq k}),
 \qquad (u,v)\longmapsto |u|.
 \label{eq:poisson_word_split_maps}
\end{align}
\end{lemma}
\begin{proof}
\leanok
Splitting a word at position $k$ and concatenating the two parts returns
the original word.  Conversely, if $uv=w$, then $|u|\leq m$,
and splitting $w$ at $|u|$ returns $u$ and $v$.  The prefix of the split
at $k\leq m$ has length $k$, so the maps
in~(\ref{eq:poisson_word_split_maps}) are inverse.
This also proves finiteness, including the unique split of the empty word.
\end{proof}
